\subsection{Additional smooth and multiplicative sites on the weighted cores}

\begin{statement}{P3F.48}{Proved}
For the actual additive cores constructed in
P3F.36, P3F.38, and P3F.41--43, P3F.46--47,
add any finite collection of logarithmic caps at smooth or
multiplicative source values, retaining their marked elliptic,
gamma, and delta collar data. Every such added cap in a sphere
member has effective order one. This holds for every allowed
cusp fan. The statement retains all order-one clutches and
assigns no parameter-return identity to them.
\end{statement}
\begin{proof}
\par\medskip\noindent\textbf{The actual source exterior.}
Let $W_0$ be the original negative-translation orbit source with
only its prescribed additive-core tubes removed. It is the
actual mapping torus of $t_{-bG}$ on that source complement,
where $b\in\{1,2,3,6\}$. Write $q=1$ or two for the number
of removed additive tubes. The complete source-cone calculations
in the cited statements give
\[
 H_2(W_0;\mathbb Z)=\mathbb Z/b\oplus\mathbb Z^q,
 \qquad H_3(W_0;\mathbb Z)=\mathbb Z^q.                         \tag{F48.1}
\]
The trivial summand is omitted for $b=1$.
An integral $H_3$ basis consists of the actual common generic-fibre
trace $F\gamma$ and $q-1$ source-boundary classes $N_i$ at
the old removed tubes. For $q=2$ their relation is
$N_1+N_2=0$. These are the marked degree-three representatives
in the actual source-cone certificates.

The source involution
\[
                         I[p,t]=[-p,-t]
\]
reverses the five-dimensional orientation, fixes the base, and
preserves every chosen tube. We verify its action on degree two
using the actual height decomposition. The rational source
lattice has the marked classes $O,F$, the root components, and
the Shioda class $\phi$. Inversion fixes $O,F$, sends
$\phi$ to $-\phi$, and takes every component to its inverse
component label. All retained multiplicative source fibres have
transitive $t_{-bG}$ component action. Their component
coinvariants therefore span only the whole-fibre direction.
The removed additive components have already been imposed by
the complement intersection map. Translation takes $\phi$ to
$\phi$ plus a nonzero rational multiple of $F$, so $F$ dies
in the rational degree-two coinvariant. Its remaining direction
is $\phi$, with negative inversion. The other $q-1$ directions
are traces of the base-annulus $H_1$ classes; source inversion
fixes those classes and time reversal negates their traces.
Thus $I=-1$ on $H_2(W_0;\mathbb Q)$.

For a prime $p\nmid b$, put $R=\mathbb Z_{(p)}$.
Equation (F48.1) makes $H_2(W_0;R)$ free. It injects into
its rationalization, so the actual action on it is exactly
$-1$. This argument does not divide by two. The full source
inversions are also retained in the $I_1^*$ certificates and
in the four cases of
\path{code/part3/iv_star_iii_joint.py}.

\par\medskip\noindent\textbf{The exact ordinary comparison.}
Remove all the added ordinary smooth and multiplicative tubes
from $W_0$, and call the resulting exterior $W$.
Let $U_j$ be these actual ordinary tubes and $E_j=\partial U_j$.
We first use their original ordinary fillings. Every
$H_2(U_j;R)$ and $H_2(E_j;R)$ is free. At a multiplicative
site its source Euler pair $(N,c)$ is primitive, since the
actual translation is transitive on components.
The ordinary cap maps on $H_2$ are surjective. Their primitive
kernels are the three smooth rims
\[
                    a\alpha,\quad b\alpha,\quad\alpha\gamma
\]
or the two actual multiplicative classes $R_j,L_j$ of P3F.40.
Denote the direct sum of these ordinary kernels by $K$.

The map $H_3(W;R)\to H_3(W_0;R)$ is onto by actual
representatives. Each old $N_i$ lies in its old additive
boundary collar, disjoint from every new tube. The common
$F\gamma$ trace can be placed over a regular source value
outside all new tubes. These representatives are the entire
integral basis in (F48.1). Hence the ordinary Mayer--Vietoris
connecting map from $H_3(W_0)$ to the new boundary $H_2$
is zero.

The next boundary map from $H_2(W_0;R)$ to
$\bigoplus H_1(E_j;R)$ is also zero. Its image lies in the
ordinary cap kernels, the primitive meridian groups
$\bigoplus R\alpha_j$. They are $I$-even, whereas its source
is $I$-odd. Equivariance makes twice each image zero.
The meridian group is torsion-free, also for $p=2$, so the
image itself is zero. The exact ordinary sequence is therefore
\[
 0\longrightarrow\bigoplus_jH_2(E_j;R)
 \longrightarrow H_2(W;R)\oplus\bigoplus_jH_2(U_j;R)
 \longrightarrow H_2(W_0;R)\longrightarrow0.
\]
Using the actual surjective cap projections gives exactly
\[
             0\longrightarrow K\longrightarrow H_2(W;R)
                  \longrightarrow H_2(W_0;R)\longrightarrow0.  \tag{F48.2}
\]
Indeed, subtract boundary lifts of all cap coordinates from an
ordinary preimage to obtain surjectivity of the last map.
A boundary preimage with zero cap coordinates lies in the
ordinary kernels, proving the middle kernel assertion.
The vanishing $H_3$ connecting map proves their injectivity.
This retains every rim and every integral index.

Let $g_-(q)$ be the Lie lift of the actual source translation
$-bG(q)$. With convention
$(p,1)\sim(t_{g_-}p,0)$, the exact local product is
\[
        (a,b,t)\longmapsto[a+b\tau(q)-t g_-(q),t].              \tag{F48.3}
\]
At $t=1$ the equality
$t_{g_-}(a+b\tau-g_-)=a+b\tau$ proves its closing.
In these coordinates $I$ is precisely $(-a,-b,-t)$ and
$\alpha$ is fixed. Thus the ordinary smooth kernels are
$I$-odd, while $ab,a\gamma,b\gamma$ are $I$-even.
On a Tate sector, write the continuation of $g_-$ as
$g_-+c_-$. The continuation is
$a\mapsto a+Nb-c_-t$, and commutes with the same minus
identity. The physical coefficient of P3F.40 is consequently
$c=c_-$, namely minus the component integer of $+bG$.
Its exact classes $z,B$ are even and $R,L$ are odd.
No involution has been extended over any replacement cap.

Both ends of (F48.2) are free and odd. Choose an $R$-linear
splitting. The involution has block matrix
\[
                           \begin{pmatrix}-I&C\\0&-I\end{pmatrix}.
\]
Its square is one, so $2C=0$. The coefficient ring has no
two-torsion, hence $C=0$. Thus all $H_2(W;R)$ is free and
odd. Every even new boundary core column maps to zero in it.
This proves the needed vanishing at $p=2$ without an
eigenprojection or an inverse of two.

\par\medskip\noindent\textbf{Eliminating the actual additive-core blocks.}
The prescribed additive-core degree-two Mayer--Vietoris map
is an integral isomorphism in each cited construction, for
all its actual collar coefficients. Project the new target
$H_2(W;R)$ plus the old cap groups onto $H_2(W_0;R)$ plus
those cap groups. On the old boundary summand its projection
is exactly that old isomorphism, by the ordinary inclusion
maps. Its image is therefore a graph over the whole projected
module and meets $K$ trivially. Eliminating that graph leaves
$K$ and the new replacement-cap targets. This is an exact
module quotient. It requires neither an involution on the old
caps nor a nonzero old Seifert determinant.

At a smooth replacement retain
\[
 \mu=m\alpha+n\gamma+\lambda_a a+\lambda_b b,
                    \qquad\gcd(m,n)=1.
\]
The cap map is the exterior square of the exact primitive
quotient by $\mu$, whose kernel is $\mu\wedge H_1(E)$.
Since the core columns vanish in $W$, its ordinary rim block is
\[
 \begin{pmatrix}
 m&0&0&-\lambda_a\\0&m&0&-\lambda_b\\0&0&m&n
 \end{pmatrix}.                                                   \tag{F48.4}
\]
Its columns are the rim projections of
$-\mu\wedge a,-\mu\wedge b,\mu\wedge\gamma,-\mu\wedge\alpha$.
For $p\mid m$, the integer $n$ is a unit in $R$.
The vector $(-\lambda_a,-\lambda_b,n)$ can therefore be
sent to $(0,0,1)$ by the actual invertible row map
\[
 \begin{pmatrix}1&0&\lambda_a/n\\0&1&\lambda_b/n\\0&0&1/n\end{pmatrix}.
\]
It preserves $mR^3$ as a submodule. The quotient of (F48.4)
is exactly $(R/mR)^2$, with both elliptic coefficients retained.

At a multiplicative replacement use the exact basis
$(z,R,B,L)$ with $z\cdot L=R\cdot B=1$ and zero
self-pairings. The exterior $z,B$ columns vanish and $R,L$
are its independent ordinary kernel. With the retained Tate
integer $j$, set $A=Nj-cn$. The actual cap rows are
\[
                 \begin{pmatrix}m&-A&s&t\\0&0&m&A\end{pmatrix}.
\]
Its isotropic-kernel equation is $mt=As$ and its actual
component transitivity gives $\gcd(m,A)=1$. Hence $s=mk$
and $t=Ak$. The actual target change $T_{\rm new}=T+kC$
gives the cap rows
\[
                 \begin{pmatrix}m&-A&0&0\\0&0&m&A\end{pmatrix}. \tag{F48.5}
\]
These are precisely the local maps of P3F.40. Their local
hypotheses are the primitive source Euler pair and
transitivity; the original global application is replaced
here by the exact zero core columns just proved.
Together with the exterior $R,L$ rows, permuting columns
to $(R,L,z,B)$ and subtracting their multiples from the last
two rows gives
$\operatorname{diag}(1,1,-m,-m)$.
The quotient is again $(R/mR)^2$.

The site kernels in (F48.2) form a direct sum, so all these
blocks occupy independent summands after the old additive
graphs are eliminated. A prime divisor of any new effective
order therefore contributes nonzero degree-two cokernel.
The complete Mayer--Vietoris sequence injects that cokernel
into $H_2$ of the closed orbit, excluding a sphere orbit.

\par\medskip\noindent\textbf{The localization applies to every sphere member.}
Keep the actual marked base meridians and gamma generator,
and quotient the orbit group by its normal elliptic generators.
Its abelianized marked Seifert relations include
\[
 m_i\alpha_i+n_i\gamma=0,\qquad
                  \sum_i\alpha_i+b_0\gamma=0,                    \tag{F48.6}
\]
where $b_0$ is the retained integer base framing.
The elliptic and delta path terms disappear under this group
quotient, not under a change of gamma marking.
If a prime divides two effective orders, their primitive norm
relations force $\gamma=0$ modulo that prime. Other
meridians are then determined, while the two indicated
meridians have only their one sum relation. The quotient has
nonzero $H_1$ over the residue field. Thus sphere members have
pairwise coprime effective orders. The $b_0$ term remains in
the argument and vanishes only after $\gamma=0$.

Every prime divisor of $b$ divides a required old good or
starred order in the actual weighted cores. For $b=2$ there
is a quadratic $I_1^*$ cap or a good $III/III^*$ order
divisible by four. For $b=3$ there is a good $IV/IV^*$ order
divisible by three. For $b=6$ both primes occur: in
$I_1^*$ plus good $IV$, or good $IV^*$ plus good $III$.
There is no restriction when $b=1$.
Consequently a prime divisor of a new smooth or multiplicative
order must be prime to $b$ in a sphere member. It is a prime
to which the preceding localized proof applies, giving a
contradiction. Every added effective order is therefore one.

The reduction to effective primitive pairs is the actual
finite-kernel quotient map of P3K.19. No ineffective factor
has merely been discarded. The finite common flop chain of
P3F.35 is the identity off the selected cusp and preserves
all additional sites, so the exclusion holds for every fan.

The certificate \path{code/part3/weighted_extra_sites.py}
exports the exact localized row and column maps, the
involution-extension identity, seven actual source-inversion
records, and the required-prime implications. Its saved
output also verifies $924$ supplementary prime cases,
including $p=2$. The ordinary extension (F48.2) is the
geometric proof supplying those algebraic blocks.
\end{proof}
